\section{Introduction}\label{sec:intro}

In 2025 the attack surface of AI systems changed shape. Two public artifacts define
the shift better than any survey: \emph{EchoLeak} (CVE-2025-32711), the first
zero-click prompt-injection exploit chain against a production AI agent, in which a
single crafted email planted instructions that Microsoft 365 Copilot later retrieved
through its own RAG pipeline and obeyed, exfiltrating organizational data through an
auto-fetched Markdown image---no user interaction required, CVSS 9.3
\cite{echolake-cve,echolake-paper}; and the MCP Inspector remote-code-execution
advisory (CVE-2025-49596), which demonstrated that the \emph{developer tooling} of
the agentic ecosystem itself---the thing used to debug every MCP server---shipped
with a critical RCE reachable through a browser
\cite{cve-2025-49596,oligo-mcp}. Alongside them, a family of MCP-STDIO command
injections across LibreChat, WeKnora, and the Cursor IDE
(CVE-2025-54136, among others) showed the same pattern recurring across vendors
\cite{ox-mcp-family}. The common structure: a system that ingests content it does
not control, grants that content reach into decisions and tools, and treats the
resulting text as trusted context.

Regulators and standards bodies moved in the same direction. The OWASP Top~10 for
LLM Applications was rebuilt for the agentic era---excessive agency climbed to
position~3 and system-prompt leakage was re-framed as \emph{hidden context exposure}
\cite{owasp2026,owasp2025}; MITRE ATLAS catalogued 101 techniques and 77
sub-techniques grounded in 68 real-world case studies, with agent-specific
tags \cite{atlas}; NIST published the Generative
AI Profile of the AI Risk Management Framework \cite{nist-ai-rmf,nist-genai}; and
the EU AI Act's high-risk obligations began to bite on a fixed statutory calendar
\cite{eu-ai-act}. The message of all four is the same: the unit of risk is no longer
\emph{a model answering a bad question}, but \emph{an agent executing a workflow
under adversary-controlled context}.

This paper is the research deliverable of a red-team engineering effort inside an
operating agentic penetration-testing framework (``red-agentic-teams''). It is not
a survey written from the outside looking in; every weakness class is mapped to
concrete modules in our own codebase, every defensive claim is checked against what
our code actually does, and every number is either computed by our scripts from
published tables or marked as unverifiable and excluded. We asked one question
throughout: \emph{of everything the 2024--2026 literature reports, what is relevant
to THIS system, in what order, and what do we change on Monday morning?}

\subsection{Contributions}\label{sec:contrib}

\begin{enumerate}[itemsep=2pt, topsep=4pt]
  \item \textbf{C1---A weakness taxonomy grounded in our operating context}
        (Section~\ref{sec:taxonomy}): fourteen classes W1--W14, each anchored to
        the OWASP 2026 list and MITRE ATLAS, and each tagged with the lens through
        which it hits us: offensive (O), defensive exposure of our own stack (D),
        or dual-use (B). The distinctive claim: for a red-team framework, the
        defensive lens is not a courtesy chapter---the framework ingests
        target-controlled content into an agent that executes tools, which makes it
        exactly the system EchoLeak attacked.
  \item \textbf{C2---An evidence base with explicit verification discipline}
        (Sections~\ref{sec:method}--\ref{sec:agentic}): the structural facts of
        direct attacks (optimization frontier 100\% jailbroken; static defenses
        failing against adaptive attackers; automation cost collapse to cents),
        the agentic evidence (indirect injection, tool/MCP compromise, memory
        poisoning, swarm cascades, misuse compliance), and a named list of claims
        we could NOT verify and therefore excluded.
  \item \textbf{C3---A relevance model with honest statistics}
        (Section~\ref{sec:stats}): $R(w) = P(w)\cdot I(w)\cdot(1-D(w))$ as a
        \emph{ranking statistic}, with Wilson intervals on every measured rate, the
        rule of three guarding ``we saw zero'' claims, Beta--Binomial pooling across
        heterogeneous studies, an at-least-one calculation over the fourteen
        classes, and a
        sensitivity analysis. We also argue---with Cox, Axelsson, and EPSS as
        witnesses---why ordinal risk matrices are the wrong instrument for this
        ranking, and why a probabilistic one is the right one.
  \item \textbf{C4---A deep mapping to our framework and a costed remediation
        queue} (Sections~\ref{sec:map}--\ref{sec:gaps}): every class mapped to the
        module that owns it today, the audit gap that remains, and nine
        recommendations ordered by relevance---the top of which (poison-resistant
        experience memory, adaptive self-probing, and a statistical layer) is
        already implemented and tested in the framework.
\end{enumerate}

\subsection{Positioning: the framework as a target}\label{sec:positioning}

Most red-teaming literature treats the AI system under test as the target and the
red team as trusted infrastructure. We hold the opposite assumption. Our framework
fetches pages from targets under test, reads their OpenAPI documents, distills
lessons into a persistent cross-run memory, and lets an LLM planner route tools by
priority. Every one of those inputs is attacker-writable in exactly the sense
EchoLeak demonstrated: a target can return \emph{any text it likes}, including text
shaped like instructions to our agent. A red-team framework that models attacks on
others but not on itself is half an instrument. This asymmetric posture---we call
it the \emph{D-lens}---drives several otherwise odd design choices in this paper:
the relevance model carries a defendedness term measured on our own code; the
recommended test plane probes our own guards with adaptive jailbreaks; and the
highest-ranked weakness class (W5, memory poisoning) is ranked by what an attacker
could do \emph{to us}, not by what we could do to a target.
